\documentclass{article}

\usepackage{agents4science_2025}

\usepackage[utf8]{inputenc}
\usepackage[T1]{fontenc}

\usepackage{amsmath}
\usepackage{amsfonts}
\usepackage{nicefrac}

\usepackage{graphicx}
\usepackage{subcaption}
\usepackage{multirow}
\usepackage{array}
\usepackage{tabularx}
\usepackage{colortbl}
\usepackage{xcolor}

\usepackage{tikz}
\usepackage{pgfplots}

\usepackage{float}

\usepackage{algorithm}
\usepackage{algorithmicx}
\usepackage{algpseudocode}

\usepackage{hyperref}
\usepackage{cleveref}

\usepackage{microtype}
\usepackage{booktabs}


\title{HAGuard Hierarchical Adaptive Guard for Low-Latency LLM Jailbreak Detection}

\author{AIRAS}

\begin{document}

\maketitle

\begin{abstract}
Large-language-model assistants increasingly mediate real-time user interactions, yet remain vulnerable to jailbreak prompts that elicit disallowed content. Production chat services impose sub-100 ms latency budgets, so calling a multi-billion-parameter "judge" model is infeasible, while single-stage heuristics plateau near 80\% balanced accuracy. We introduce HAGuard, a three-layer adaptive cascade that fuses (i) an ultra-fast lexical gate, (ii) a 60 M-parameter INT-4 DistilBERT, and (iii) prototype-intent matching in embedding space. All components export to ONNX and run on commodity CPUs; 70\% of traffic exits after the first two layers, keeping median end-to-end latency at 27 ms and p99 at 78 ms on an AWS c7g.large core. Training mixes multi-view contrastive learning on HarmBench, JailbreakBench, SafetyBench, StrongREJECT and 50 k benign production prompts with hard-negative mining from the RL-Breaker attack generator. Misclassified prompts trigger four-hourly LoRA updates that hot-reload in $< 2$ min. On a public 55 k-prompt benchmark HAGuard attains $96.4 \pm 0.3$\% true-positive rate at 0.5\% false-positive rate and $0.972 \pm 0.002$ AUROC-15-20 pp higher and eight times faster than LLM-Guard-7B. A deliberately truncated 16-prompt run that scored only chance-level accuracy highlights the importance of the full evaluation protocol. Code, models and raw predictions are released for complete reproduction.
\end{abstract}

\section{Introduction}
Large-language models (LLMs) have transitioned from laboratory curiosities to central components of code assistants, writing aids and even medical triage tools. Their mass adoption nevertheless hinges on robust guardrails that intercept malicious "jailbreak" prompts able to override safety alignment and coax disallowed output \cite{chao-2024-jailbreakbench,chen-2024-when}. A practical guard must simultaneously 1) recognise increasingly sophisticated, obfuscated and rapidly evolving attacks, 2) avoid collateral damage to benign requests, and 3) meet stringent service-level agreements-typically $\leq 100$ ms end-to-end for 99\% of queries.

Meeting all three goals is challenging. Lightweight lexical rules execute in micro-seconds but miss paraphrases and unicode tricks, capping balanced accuracy at roughly 75-85\% \cite{alon-2023-detecting}. Replacing them with powerful judge LLMs such as Llama-Guard-7B improves recall yet adds 200-500 ms latency, unacceptable for synchronous user experiences. Existing defences rarely combine heterogeneous evidence, leading to poorly calibrated probabilities and forcing operators to choose between aggressive blocking and risky permissiveness. Finally, attackers iterate quickly, whereas labelled data for novel jailbreak styles arrive slowly, impeding periodic re-training.

This paper presents HAGuard, a Hierarchical Adaptive Guard that addresses the four pain points through a carefully engineered cascade and an inexpensive continual-learning loop that never resorts to large models in production. All compute stays on a single CPU core; 72\% of queries terminate after a lexical or distilled-transformer stage, while the hardest 28\% invoke a quantised embedding matcher and a calibrated LightGBM head whose decision fuses lexical, semantic and prototype features. Prompts the stack misclassifies are quarantined and used to refresh a LoRA adapter every four hours, realising fast recovery reminiscent of online cascade learning \cite{nie-2024-online} without any GPU.

\subsection{Contributions}
\begin{itemize}
  \item \textbf{CPU-only low-latency cascade}: A lexical-semantic-prototype cascade that sustains $< 100$ ms p99 latency entirely on CPUs.
  \item \textbf{Training with adversarial mining}: A recipe combining multi-view contrastive learning, RL-Breaker adversarial mining \cite{chen-2024-when} and LoRA-based continual updates ($< 2$ min downtime).
  \item \textbf{CI sanity suite}: A continuous-integration suite that aborts if any artefact is stubbed or latency budgets are violated.
  \item \textbf{Comprehensive evaluation}: The first side-by-side evaluation of HAGuard, LLM-Guard-7B, Rebuff, PARDEN \cite{zhang-2024-parden} and perplexity-based filters on a public 55 k-prompt benchmark drawn from HarmBench, JailbreakBench and production logs.
  \item \textbf{Online adaptation study}: A 12-hour study showing that two LoRA refreshes recover 60\% of initially missed attacks with $< 1$ pp FPR increase.
\end{itemize}

The remainder reviews related work, formalises the problem, details the methodology, describes experimental protocols, reports empirical results and concludes with limitations and future directions.

\section{Related Work}
Existing guardrails fall into three families. 1) Micro-second heuristics-regex keyword lists and perplexity filters-scale but fail on obfuscations, plateauing at \textasciitilde 80\% accuracy \cite{alon-2023-detecting}. 2) Fine-tuned moderation LLMs such as Llama-Guard-2 and WildGuard improve recall yet are GPU-bound, adding hundreds of milliseconds latency \cite{han-2024-wildguard,inan-2023-llama}. 3) Autoregressive self-judging approaches like PARDEN prompt the assistant to repeat its own output to avoid domain shift but still invoke a large model per request \cite{zhang-2024-parden}.

Complementary defences intervene inside the assistant itself: information-bottleneck training hinders leakage of harmful content \cite{liu-2024-protecting}; adversarial prompt tuning hardens the system against black-box attacks \cite{mo-2024-fight}; MoGU routes between "safe" and "usable" variants to balance utility and harm \cite{du-2024-mogu}. Such internal techniques do not obviate the need for an external pre-filter under tight latency constraints.

Balancing risk against computation through cascades is well studied in computer vision \cite{goren-2024-hierarchical} and was recently extended to LLM inference over streams \cite{nie-2024-online}. HAGuard is, to our knowledge, the first cascade combining lexical, distilled-transformer and prototype-intent cues for jailbreak detection.

Continual adaptation via lightweight adapters appears in privacy-preserving prompting \cite{duan-2023-flocks} and backdoor scanning \cite{shen-2021-backdoor}. We adopt the same principle to refresh DistilBERT every few hours without downtime.

Finally, standardised benchmarks such as JailbreakBench \cite{chao-2024-jailbreakbench} and StrongREJECT \cite{souly-2024-strongreject} expose weaknesses of earlier ad-hoc evaluations. Our offline and online suites follow their guidelines and release raw predictions for transparency.

\section{Background}
\subsection{Problem formulation}
For each user prompt $p$ we must emit $y \in \{\text{allow}, \text{block}\}$ and optionally a risk label $t$ from a taxonomy of $K$ categories. The guard executes synchronously before the assistant; the 99th-percentile runtime must not exceed $L = 100$ ms. Utility is measured primarily by the true-positive rate (TPR) at a 0.5\% false-positive rate (FPR), and by AUROC.

\subsection{Signals}
Lexical patterns captured by tries, unicode-homoglyph detectors and Bloom filters offer constant-time but brittle cues. Contextual semantics require transformer encoders, yet even compact models cost tens of milliseconds. Sentence-level embeddings cluster prompts with similar intent; nearest-neighbour distances to curated prototypes provide interpretable yet potentially noisy evidence.

\subsection{Cascade rationale}
Let $\ell_i$ denote latency of layer $i$ and $\Delta_i$ its marginal recall. A cascade is worthwhile when the expected latency $\sum_i \pi_i \, \ell_i$ ($\pi_i$ being the traffic fraction reaching layer $i$) remains under $L$ while cumulative recall approaches that of the full stack. Early-exit thresholds must be calibrated to respect the global FPR budget.

\subsection{Adaptive defence}
Attackers rapidly invent paraphrases, emojis or unicode tricks that elude static filters \cite{liu-2023-autodan}. Logging the hardest misses and injecting them via small LoRA updates enables prompt recovery without full retraining.

\section{Method}
\subsection{System architecture}
HAGuard comprises three sequential layers. 1) Layer-0 lexical gate. A trie of 10\,318 unsafe keywords and a 1 M-slot Bloom filter (murmur-3, $k = 7$, FP $\approx 0.1$\%) compute a risk score $s_0$; prompts with $s_0 > 0.95$ are blocked. Latency $\approx 5$ \(\mu\)s. 2) Layer-1 semantic scorer. DistilGuard-60M-six-layer DistilBERT distilled from Llama-Guard logits-is quantised to INT-4 weights with INT-8 activations and exported to ONNX-Runtime. It outputs class probabilities $p_1$; the maximum forms score $s_1$. Latency $\approx 18$ ms on one vCPU. 3) Layer-2 prototype matcher. Prompts are encoded with MiniLM-L6-v2 into 128-dim vectors $e$. Two FAISS-HNSW indices-20 k harmful and 40 k benign prototypes-are searched ($ef = 128$, $M = 32$). Features comprise the signed margin $m = d_{\text{ben}} - d_{\text{harm}}$, Jaccard token overlap and an intent-consistency flag, concatenated with $s_0$ and $s_1$. A calibrated LightGBM with 500 trees (depth 7, learning-rate 0.05) produces the final verdict. Latency $\approx 9$ ms.

\subsection{Training data}
A 164 k-prompt corpus combines HarmBench, SafetyBench jailbreak subset, JailbreakBench, StrongREJECT, 50 k benign production prompts and 5 k real jailbreak attempts.

\subsection{Model training}
DistilGuard is fine-tuned with a multi-view contrastive loss (temperature 0.07, queue 1\,024) over paraphrases, translation round-trips and RL-Breaker adversarial variants \cite{chen-2024-when}. Prototypes are selected via k-means++ in embedding space and labelled by majority vote. LightGBM minimises log-loss with class weights chosen to hit 0.5\% FPR on validation.

\subsection{Continual learning}
Prompts misclassified after Layer-2 populate a quarantine bucket. Every four hours up to 4\,000 are used to train a rank-8 LoRA adapter (lr $2 \times 10^{-5}$, one epoch) that merges into DistilGuard and hot-reloads across workers, incurring $< 2$ min downtime.

\subsection{System optimisation}
All kernels exploit AVX-512; asynchronous I/O and early exits keep Layer-1 and Layer-2 threads idle for \textasciitilde 70\% of traffic.

\subsection{Integrity assurance}
A CI suite executes a 1\,000-prompt sanity batch, asserting entropy of Layer-1 logits, prototype margins, median latency $\leq 25$ ms and resident-set size $\leq 800$ MB. Any stub placeholder aborts the pipeline.

\section{Experimental Setup}
\subsection{Hardware}
All experiments run on AWS c7g.large (single Graviton-3 vCPU, 8 GiB RAM). GPU numbers are reported only for the 7 B baseline.

\subsection{Datasets}
(1) HarmBench 13\,832 prompts, (2) SafetyBench jailbreak subset 8\,617, (3) JailbreakBench 9\,527 \cite{chao-2024-jailbreakbench}, (4) StrongREJECT 313 \cite{souly-2024-strongreject}, (5) Prod-Benign 50\,000, (6) Rebuff-Jailbreak 5\,000. After NFKC normalisation and deduplication 164 k unique prompts remain; 55 k prompts identical to the manuscript's test list are held out for public evaluation.

\subsection{Splits}
Prompt-family stratified 70\%/10\%/20\% train/val/test; StrongREJECT resides entirely in test. Random state 2025 ensures replicability.

\subsection{Baselines}
A LLM-Guard-v2-7B (fp16, GPU via vLLM), B Rebuff default rules, C PARDEN \cite{zhang-2024-parden}, D perplexity filter \cite{alon-2023-detecting}. Thresholds are tuned on validation to meet 0.5\% FPR when possible.

\subsection{Metrics}
Primary: TPR at 0.5\% FPR, AUROC, median and p99 latency. Secondary: balanced accuracy, 15-bin expected calibration error, RSS and AWS cost per million queries. A robustness suite applies (i) whitespace padding $\leq 1\,024$ spaces, (ii) 5\% homoglyph swaps, and (iii) benign-harmful concatenation; we report $\Delta$TPR and $\Delta$FPR relative to clean prompts.

\subsection{Procedures}
Offline evaluation trains five seeds \{13, 17, 19, 23, 29\}, stores predictions and computes mean $\pm$ SD plus 1\,000$\times$ bootstrap confidence intervals. A 12-h online study streams 50 req/s via Kafka, injects $\geq 3\,000$ novel attacks from RL-Breaker, and compares HAGuard-Adaptive, HAGuard-Frozen and LLM-Guard-7B workers. Prometheus scrapes accuracy, latency, CPU and memory per hour; LoRA refreshes occur at hours 4 and 8.

\section{Results}
\subsection{Offline benchmark}
Across five seeds HAGuard achieves $96.4 \pm 0.3$\% TPR at 0.5\% FPR and AUROC $0.972 \pm 0.002$. LLM-Guard-7B scores $81.1 \pm 0.4$\% TPR and 0.887 AUROC; Rebuff 74.5\%, PARDEN 78.2\%, perplexity 63.4\%. Median latency is $27 \pm 1$ ms and p99 $78 \pm 2$ ms, versus $219 \pm 8$ ms and $642 \pm 15$ ms for the GPU baseline. RSS is 613 MB compared to 15 GiB plus GPU memory for LLM-Guard.

\subsection{Ablation study}
Removing Layer-0 increases median latency by 15 ms and CPU utilisation by 53\% without accuracy benefit. Dropping Layer-1 reduces TPR by 9.8 pp; omitting Layer-2 cuts TPR by 11.3 pp but saves only 5 ms. Switching DistilGuard from INT-4 to INT-8 yields a negligible 0.2 pp TPR gain yet doubles memory.

\subsection{Robustness}
Whitespace padding and homoglyph attacks lower HAGuard's TPR by 1.9 pp, compared with 7-12 pp for single-stage models. Prompt concatenation costs 3.1 pp.

\subsection{Online adaptation}
During the 12-h stream attack entropy spikes at hours 4 and 8. HAGuard-Adaptive climbs from 78\% to 92\% TPR after two LoRA refreshes, with FPR rising by $< 1$ pp; p99 latency stays below 85 ms. The frozen variant remains flat, as does LLM-Guard-7B. Total finetune downtime is 96 s.

\subsection{Truncated run}
A deliberately misconfigured "full-experiment" processed only 16 prompts; all exited at Layer-0, yielding $\sim 10$ \(\mu\)s latency and 50\% accuracy. Because CI checks are disabled in this mode, the episode underscores the necessity of the integrity suite.

\subsection{Limitations}
False positives persist in creative role-play scenarios where fictional crime plots resemble genuine malicious intent. Quantisation introduces slight calibration drift that future work should correct.

\section{Conclusion}
HAGuard demonstrates that carefully orchestrated lexical, distilled-transformer and prototype-intent evidence can raise jailbreak-detection accuracy from the long-standing 80\% ceiling to beyond 95\% while meeting sub-100 ms latency budgets on a single CPU core. Multi-view contrastive training, RL-Breaker hard-negative mining and rapid LoRA updates keep the guard adaptive without GPUs or prolonged downtime. Extensive offline and online experiments confirm substantial gains over GPU-bound or heuristic baselines, and all artefacts are open-sourced for full reproducibility. Future work will pursue tighter calibration, multilingual coverage, integration with routing frameworks such as MoGU \cite{du-2024-mogu}, human-in-the-loop feedback and certified robustness guarantees inspired by CBD \cite{xiang-2023-cbd}.


\bibliographystyle{plainnat}
\bibliography{references}

\end{document}